\subsubsection{五阶递推特征多项式的数域结构与分裂型指纹}\label{subsec:pom-s5-galois-arithmetic}

令五阶碰撞矩递推的特征多项式为
\[
P_5(\lambda)
:=\lambda^5-2\lambda^4-11\lambda^3-8\lambda^2-20\lambda+10\in\ZZ[\lambda].
\]
记其五个根为 \(\alpha_1,\dots,\alpha_5\)，并记分裂域为 \(L_5\)。

\begin{proposition}[五阶特征多项式的满对称伽罗瓦群]\label{prop:pom-s5-galois-s5}
有同构
\[
\mathrm{Gal}(L_5/\QQ)\cong S_5.
\]
并且
\[
\Disc(P_5)=-16107783120=-2^{4}\cdot 3^{4}\cdot 5\cdot 11\cdot 13\cdot 17383,
\]
从而 \(\Disc(P_5)\notin \QQ^{\times 2}\)。
\leanverified{paper_pom_s5_galois_s5_seeds}
\leanverified{paper_pom_s5_galois_s5_package}
\leanverified{paper_pom_s5_galois_s5}
\end{proposition}
\begin{proof}
取素数 \(p=17\)。在 \(\FF_{17}[\lambda]\) 上，\(P_5\bmod 17\) 不可约，从而 \(P_5\) 在 \(\QQ[\lambda]\) 上不可约；因此伽罗瓦群 \(G:=\Gal(L_5/\QQ)\le S_5\) 作为置换群在五个根上是传递的。
由于 \(17\nmid \Disc(P_5)\)，故 \(p=17\) 在 \(L_5\) 中不分歧；由 Dedekind--Frobenius 对应可知，\(\bmod 17\) 不可约蕴含某个 Frobenius 元在 \(G\) 中为 \(5\)-循环。
\par
再取素数 \(p=29\)。同样 \(29\nmid \Disc(P_5)\) 不分歧，并且
\[
P_5(\lambda)\equiv (\lambda+4)(\lambda+8)\bigl(\lambda^3-14\lambda^2+9\lambda+13\bigr)\pmod{29},
\]
故 \(G\) 含循环型为 \((3+1+1)\) 的元素，特别地含一个 \(3\)-循环。
\par
含 \(5\)-循环的传递子群只能为 \(C_5,D_5,\mathrm{AGL}(1,5),A_5,S_5\)；其中前三者不含 \(3\)-循环。
因此 \(G\in\{A_5,S_5\}\)。
而 \(\Disc(P_5)\notin\QQ^{\times2}\) 等价于符号表示非平凡，从而 \(G\not\subseteq A_5\)。
于是 \(G=S_5\)。证毕。
\end{proof}

\begin{proposition}[Perron 数域的单生成元性与精确判别式]\label{prop:pom-s5-field-discriminant-monogenic}
令 \(\rho_5\) 为 \(P_5\) 的 Perron 根，并记
\(
K_5:=\QQ(\rho_5)
\)。
则 \(\mathcal O_{K_5}=\ZZ[\rho_5]\)，并且
\[
\Disc(K_5)=\Disc(P_5)=-16107783120.
\]
特别地，\(K_5\) 的分歧素数集合恰为 \(\{2,3,5,11,13,17383\}\)，且判别式平方类满足
\[
\Disc(K_5)\bmod \QQ^{\times 2}=-12428845.
\]
\end{proposition}
\begin{proof}
对 \(\ZZ[\rho_5]\) 应用极大整环的 round-two 算法可得 \([\mathcal O_{K_5}:\ZZ[\rho_5]]^2=\Disc(P_5)/\Disc(K_5)=1\)，从而 \(\mathcal O_{K_5}=\ZZ[\rho_5]\) 且 \(\Disc(K_5)=\Disc(P_5)\)。
分歧素数集合与平方类断言由判别式分解与去平方因子直接读出。
可复现脚本：\texttt{scripts/exp\_fold\_collision\_moment\_charpoly\_galois\_audit\_q5.py}。证毕。
\end{proof}
\leanverified{paper\_pom\_s5\_field\_discriminant\_monogenic}

\begin{proposition}[由二根对称构造得到的指数 \(10\) 子域与判别式平方类刚性]\label{prop:pom-s5-two-subset-degree10}
对任意 \(1\le i<j\le 5\)，置
\[
\sigma_{ij}:=\alpha_i+\alpha_j,\qquad \pi_{ij}:=\alpha_i\alpha_j.
\]
则 \(\sigma_{ij}\) 与 \(\pi_{ij}\) 的最小多项式次数均为 \(10\)，并分别满足不可约十次多项式
\[
R^{(+)}_5(s)
=s^{10}-8s^9-9s^8+158s^7+107s^6-1482s^5-711s^4+2262s^3+432s^2+10350s+7200,
\]
\[
R^{(\times)}_5(p)
=p^{10}+11p^9+36p^8+436p^7-1540p^6+7080p^5-18700p^4-2900p^3+24000p^2+8000p+10000.
\]
并且在 \(2\)-子集上的忠实传递表示下有
\[
\Gal\!\left(R^{(+)}_5/\QQ\right)\cong S_5,
\qquad
\Gal\!\left(R^{(\times)}_5/\QQ\right)\cong S_5.
\]
此外两者的判别式满足同一平方类刚性
\[
\Disc\!\left(R^{(+)}_5\right)\equiv \Disc\!\left(R^{(\times)}_5\right)\equiv \Disc(P_5)^3
\pmod{\QQ^{\times 2}},
\]
且其显式分解为
\[
\Disc\!\left(R^{(+)}_5\right)
=-2^{18}\cdot 3^{18}\cdot 5^{7}\cdot 11^{3}\cdot 13^{3}\cdot 17383^{3}\cdot 1193^{2}\cdot 2237^{2},
\]
\[
\Disc\!\left(R^{(\times)}_5\right)
=-2^{32}\cdot 3^{20}\cdot 5^{23}\cdot 11^{3}\cdot 13^{3}\cdot 17383^{3}\cdot 71^{2}\cdot 79^{2}.
\]
\leanverified{paper\_pom\_s5\_two\_subset\_degree10\_seeds}
\leanverified{paper\_pom\_s5\_two\_subset\_degree10\_package}
\leanverified{paper\_pom\_s5\_two\_subset\_degree10}
\end{proposition}
\begin{proof}
令 \(s=x+y\)。对方程 \(P_5(x)=0\)、\(P_5(y)=0\) 与 \(s-(x+y)=0\) 作消元，得到结果式分解
\[
\Res_x\bigl(P_5(x),P_5(s-x)\bigr)=\bigl(32\cdot P_5(s/2)\bigr)\cdot \bigl(R^{(+)}_5(s)\bigr)^2.
\]
其中因子 \(32\cdot P_5(s/2)\) 对应对角解 \(x=y\)（即 \(s=2x\)），而平方出现来自 \((x,y)\) 与 \((y,x)\) 给出同一和。
故对任意 \(i\neq j\)，取 \((x,y)=(\alpha_i,\alpha_j)\) 得 \(R^{(+)}_5(\sigma_{ij})=0\)。
\par
同理令 \(p=xy\)。用约束 \(p-xy=0\) 消去并取结果式，可得某个对角因子 \(\widetilde P(p)\)（对应 \(x=y\)）与平方因子 \(\bigl(R^{(\times)}_5(p)\bigr)^2\)，从而 \(R^{(\times)}_5(\pi_{ij})=0\)（\(i\neq j\)）。
\par
由命题 \ref{prop:pom-s5-galois-s5}，\(G=\Gal(L_5/\QQ)\cong S_5\)。
而 \(S_5\) 在 \(2\)-子集集合 \(\{\{i,j\}:1\le i<j\le 5\}\) 上传递且忠实，故 \(\sigma_{ij}\) 与 \(\pi_{ij}\) 的共轭轨道大小均为 \(10\)，从而其最小多项式次数为 \(10\)，并且相应十次多项式不可约；其伽罗瓦群为该 \(10\) 点表示下的 \(S_5\) 之像，仍同构于 \(S_5\)。
\par
判别式的显式分解与平方类刚性均可由上述多项式直接计算并审计；可复现脚本：\texttt{scripts/exp\_fold\_collision\_moment\_charpoly\_galois\_audit\_q5.py}。证毕。
\end{proof}

\begin{proposition}[五阶比值型元素的指数 \(20\) 子域与忠实 \(20\) 点表示]\label{prop:pom-s5-ordered-ratio-degree20}
\leanverified{paper\_pom\_s5\_ordered\_ratio\_degree20}
对有序对 \(i\neq j\) 定义
\[
u_{ij}:=\frac{\alpha_j^2}{\alpha_i}.
\]
则 \(u_{ij}\) 在 \(\QQ\) 上的最小多项式次数为 \(20\)，并满足不可约 \(20\) 次多项式
\[
\begin{aligned}
\mathfrak N_5(u)=\;&500u^{20}-25000u^{19}-177700u^{18}-3823100u^{17}+40765660u^{16}\\
&+1496266360u^{15}+9983236609u^{14}+36022632766u^{13}+64996856413u^{12}\\
&-87017173082u^{11}-763519868228u^{10}+2321405618236u^9-2486082754184u^8\\
&+244212607720u^7+786872105380u^6-96473453600u^5-197869496000u^4\\
&-17361320000u^3-2353600000u^2-46000000u+5000000,
\end{aligned}
\]
并且
\[
\Gal(\mathfrak N_5/\QQ)\cong S_5,
\]
其中同构由 \(S_5\) 在有序对集合 \(\{(i,j):i\neq j\}\) 上的忠实 \(20\) 点作用给出。
\end{proposition}
\begin{proof}
令 \(u\) 为参数，考虑代数簇
\[
P_5(x)=0,\qquad P_5(y)=0,\qquad ux-y^2=0.
\]
先对 \(y\) 消元得
\[
\Res_y\bigl(P_5(y),ux-y^2\bigr)=u^5x^5-26u^4x^4+49u^3x^3+416u^2x^2+560ux-100.
\]
再对 \(x\) 消元，得到
\[
\Res_x\Bigl(P_5(x),\Res_y(P_5(y),ux-y^2)\Bigr)
=-200\cdot P_5(u)\cdot \mathfrak N_5(u).
\]
其中因子 \(P_5(u)\) 对应对角解 \(y=x\)（此时 \(ux=x^2\Rightarrow u=x\)），而其余因子即为 \(\mathfrak N_5\)。
因此对 \((x,y)=(\alpha_i,\alpha_j)\) 且 \(i\neq j\) 有 \(\mathfrak N_5(u_{ij})=0\)。
\par
由命题 \ref{prop:pom-s5-galois-s5}，\(S_5\) 在有序对集合 \(\{(i,j):i\neq j\}\) 上 \(2\)-传递，故任一 \(u_{ij}\) 的共轭轨道大小为 \(20\)，从而其最小多项式次数为 \(20\)，并推出 \(\mathfrak N_5\) 不可约。
忠实性可由 \(2\)-传递性推出，因此 \(\Gal(\mathfrak N_5/\QQ)\) 与该 \(20\) 点表示下的 \(S_5\) 同构。
可复现脚本：\texttt{scripts/exp\_fold\_collision\_moment\_charpoly\_galois\_audit\_q5.py}。证毕。
\end{proof}

\begin{corollary}[由 Čebotarev 定理导出的模素分解型密度律]\label{cor:pom-s5-chebotarev-splitting-density}
对任意不分歧素数
\(
p\nmid 2\cdot 3\cdot 5\cdot 11\cdot 13\cdot 17383
\)，
\(P_5(\lambda)\bmod p\) 的分解型在素数集合中按自然密度分布，并等于 \(S_5\) 中相应循环型共轭类的比例。具体为：
\begin{itemize}
  \item 不可约五次型（循环型 \((5)\)）：密度 \(1/5\)；
  \item \((4+1)\) 型：密度 \(1/4\)；
  \item \((3+2)\) 型：密度 \(1/6\)；
  \item \((3+1+1)\) 型：密度 \(1/6\)；
  \item \((2+2+1)\) 型：密度 \(1/8\)；
  \item \((2+1+1+1)\) 型：密度 \(1/12\)；
  \item 全分裂型：密度 \(1/120\)。
\end{itemize}
\leanverified{paper\_pom\_s5\_chebotarev\_splitting\_density}
\end{corollary}
\begin{proof}
由命题 \ref{prop:pom-s5-galois-s5}，\(L_5/\QQ\) 为 \(S_5\)-伽罗瓦扩张。
对不分歧素数 \(p\)，Frobenius 共轭类的循环型与 \(P_5\bmod p\) 的不可约因子次数分拆一致。
由 Čebotarev 密度定理，给定共轭类的密度等于其在 \(S_5\) 中的比例；逐一计数各循环型在 \(S_5\) 中的元素数即得。证毕。
\end{proof}

\begin{conjecture}[碰撞矩特征域与比值域的满对称性]\label{conj:pom-moment-charpoly-ratio-full-symmetric}
对本文由经审计整系数递推给出的每个 \(q\ge 4\)，记 \(P_q\in\ZZ[\lambda]\) 为最小递推的特征多项式，次数为 \(d_q\)，根为 \(\alpha_1,\dots,\alpha_{d_q}\)。
猜想：
\begin{enumerate}
  \item \(P_q\) 在 \(\QQ\) 上不可约且 \(\Gal(P_q/\QQ)\cong S_{d_q}\)；
  \item 对任意 \(i\neq j\)，比值型元素 \(u_{ij}:=\alpha_j^2/\alpha_i\) 的共轭轨道大小为 \(d_q(d_q-1)\)，从而其最小多项式不可约且其伽罗瓦群为 \(S_{d_q}\) 在有序对上的忠实表示之像；
  \item 若由自然子表示（例如 \(2\)-子集表示与有序对表示）构造相应子域，则其判别式平方类在 \(\QQ^\times/(\QQ^\times)^2\) 中由 \(\Disc(P_q)\) 的幂所刚性决定。
\end{enumerate}
\end{conjecture}
